\documentclass[journal]{IEEEtran}
% General project reference: no six-page limit.
% Short EIG conference draft: ../TPEC_conference/eig_power_grid_conference.tex.
\usepackage{amsmath,amssymb,booktabs,graphicx,xcolor}
\usepackage[hidelinks]{hyperref}

\newcommand{\todo}[1]{\textcolor{blue}{[TODO: #1]}}
\newcommand{\E}{\mathbb{E}}
\newcommand{\mocu}{\mathrm{MOCU}}
\newcommand{\ocu}{\mathrm{OCU}}

\title{Objective-Driven Sequential Bayesian Experimental Design for
Power-Grid Dynamic Learning and Control}
\author{\todo{Authors and affiliations}%
\thanks{\todo{Funding and corresponding-author information.}}}

\begin{document}
\maketitle

\begin{abstract}
This document records the formulation, methods and evaluation protocols of
the objective-driven sequential Bayesian experimental-design project.
Non-reset power-grid probing couples parameter learning to the physical
state from which subsequent experiments and terminal control operate.
The framework includes expected information gain (EIG), minimum safe control
(MSC), and finite-loss mean objective cost of uncertainty (MOCU), with
separate treatment of posterior objectives, physical safety and numerical
approximations. DAD, RL-sBOED, Step-DAD, Myopic, Fixed and Random provide the
core method comparison; MoE-sBOED is an optional extension. IEEE9 and the
reduced IEEE14 EIG backend, the separate SIR benchmark, and the pending
IEEE30 extension have distinct implementation and validation scopes.
Current conference settings are documented as one campaign within this
broader project, not as a restriction on its scientific content. Empirical
claims require completed and verified results under the stated protocol.
\end{abstract}
\begin{IEEEkeywords}
Bayesian experimental design, power-system identification, adaptive probing,
expected information gain, minimum safe control, MOCU, mixture of experts.
\end{IEEEkeywords}

\section{Introduction}
\label{sec:introduction}
\paragraph{Document scope}
This is the general project reference, retaining the EIG, MSC and MOCU
formulations, physical models, numerical methods, benchmark status and
optional MoE-sBOED description. It is not restricted to six pages.
The separate source \texttt{TPEC\_conference/eig\_power\_grid\_conference.tex} is the short
IEEE conference draft for the six-method IEEE9/IEEE14 EIG study.
The conference excludes MoE-sBOED experiments, but this general document
retains its optional formulation in Section~\ref{sec:optional-moe}.
Control-objective sections remain part of the project; their presence does
not imply that control-safety metrics were computed by an EIG-only campaign.

Uncertain inertia and effective frequency response influence both how a power
system reacts to an injected probe and how much supplementary power is needed
following a disturbance. Controlled excitation provides a way to learn these
quantities rather than waiting for an uncontrolled event. Peng et
al.~\cite{peng2024probing} develop a probing-based approach to joint inertia
and frequency-response estimation. This motivates the present identification
setting, while the reduced multi-machine dynamics, prior distributions,
observation summary and terminal-control task used here are stated study
assumptions rather than consequences of that estimation method.

The experimental-design question is which excitation to apply next, given the
measurements already collected. A probe is useful when its response helps
resolve uncertainty relevant to subsequent decisions, but information about
parameters and improvement in an engineering objective need not coincide.
Two posterior distributions can have different parameter uncertainty while
inducing the same control decision. Conversely, a relatively small change in
posterior mass near a safety threshold can change the selected control.
These observations motivate comparing parameter-oriented expected information
gain (EIG) with objectives defined directly through the terminal decision.
The objective-based view follows the distinction between model uncertainty
and its operational cost in MOCU~\cite{yoon2013mocu,boluki2018generalized}.

This work considers a finite sequence of duration-controlled probes followed
by a continuous supplementary power injection. Probe bus and amplitude are
fixed, and each probe is observed over a fixed recording window. The physical
system is initialized once, before the first probe, and its state carries
between windows. Consequently, changing a duration affects the current
measurement, the state entering later probes, and the state from which
terminal control begins. A likelihood evaluated from equilibrium after every
probe would represent a different experiment. Similarly, reordering an
observed history would generally change its physical interpretation.

The conference observation is one noisy signed endpoint
rate-of-change-of-frequency (RoCoF) value per 3.5-second window, computed from
the final two frequency samples at 40 Hz. It becomes available at window end.
The maximum-absolute-RoCoF sensor is retained as a separate historical variant;
its results and peak-switch derivative requirements are not interchangeable
with the endpoint protocol. Numerical differentiation must include both the
observation and the carried physical state.

Three objectives connect the same probing model to different notions of
experimental value. EIG rewards information about the uncertain parameters.
Minimum safe control (MSC) minimizes the expected magnitude selected by a
posterior-coverage controller. Finite-loss MOCU minimizes the posterior excess
loss of a common control relative to controls selected with knowledge of the
model. MSC and MOCU can select the same control for an individual posterior
under the stated assumptions, yet assign different values to probe policies.
Physical safety is evaluated separately on the true simulated system; a
posterior coverage level is not identified with an empirical safety rate.

\subsection{Research Gap}
The unresolved issue addressed here is the joint effect of objective choice
and policy adaptation in an experiment whose probes also change the future
physical state. A method comparison based only on parameter information does
not determine its control consequences. A comparison that resets the simulator
between probes does not represent the stated operating sequence. These two
issues must be treated together, while keeping optimization error and
posterior approximation error distinct from differences between objectives.
The literature discussed next supplies the identification, decision-theoretic
and policy-learning foundations; the present study specifies their interaction
within one reproducible non-reset power-grid benchmark.

\subsection{Research Purpose and Contributions}
The research purpose is to determine how the definition of experimental value
changes sequential probe selection when learning and physical state evolution
are coupled. The study asks three questions. First, how do information-oriented
and control-oriented objectives value the same admissible histories? Second,
what additional value do observation dependence, horizon planning and online
policy refinement provide under matched computational settings? Third, which
numerical approximations materially affect the reliability of that comparison?
These are empirical questions within the specified model, not assumptions
that an adaptive method must outperform a baseline.

The contributions are threefold. We formulate an ordered non-reset probing
problem that connects uncertain swing dynamics to a posterior-dependent
terminal action. We specify EIG, MSC and finite-loss MOCU within the same
physical and observation model while separating them from true-system safety.
We develop a reproducible comparison of six design strategies, including
state-carrying inference, feasible continuous actions, objective-specific
optimization and numerical verification. The contribution is the application
formulation and carefully specified comparison; the policy families themselves
are adapted from the cited literature. IEEE9 and the reduced five-machine IEEE14 model are implemented EIG benchmarks,
and the scope of any empirical conclusion is limited by the completed
experiments, seed coverage and model assumptions.

\subsection{Organization of the Paper}
Section~\ref{sec:related-work} positions the study within the literature.
Section~\ref{sec:formulation} defines the physical model, sequential experiment
and design objectives. Section~\ref{sec:methodology} develops the six solution
methods and their numerical implementation. Section~\ref{sec:experimental-protocol}
specifies the experimental setup and evaluation procedures.
Sections~\ref{sec:results} and~\ref{sec:conclusion} are reserved for the results,
discussion and conclusions of the completed study.

\section{Related Work}
\label{sec:related-work}
\subsection{Probing for Dynamic Parameter Identification}
Probing-based identification uses a designed excitation to expose dynamic
responses that are difficult to distinguish from passive measurements alone.
Peng et al.~\cite{peng2024probing} jointly estimate inertia and frequency
response using an injected probing signal, with a regression-based estimator
and power-hardware-in-the-loop validation. Their work motivates learning the
two effects together. Our question concerns selection of successive probe
durations under posterior uncertainty and a downstream control objective.
The scalar endpoint-RoCoF observation, state-carry model and Bayesian prior here
are additional modeling choices; their validity is not established by the
hardware validation of a different estimator or measurement model.

\subsection{Amortized and Semi-Amortized Bayesian Design}
DAD learns a policy before deployment so that the next experimental design can
be obtained from the accumulated observations by a network evaluation.
Foster et al.~\cite{foster2021dad} introduce sequential contrastive information
bounds for this purpose. We use the sequential prior-contrastive objective
for EIG, while replacing exchangeable history representations with ordered
stage slots appropriate to a state-carrying simulator. Optimizing MSC or MOCU
with a deterministic policy is an objective-specific adaptation, not an
assertion that the original DAD training objective already represents those
control criteria.

Blau et al.~\cite{blau2022rlsboed} formulate sequential experimental design as
a reinforcement-learning problem, allowing policy optimization without
requiring derivatives through the probabilistic simulator. Our RL-sBOED
implementation uses a continuous REDQ-style actor--critic and rewards formed
from changes in the prefix objective. Stepwise rewards do not make this a
myopic method: the critic accounts for later rewards over the remaining
horizon. Both its entropy regularization and its numerical budgets must be
distinguished from the information score reported at evaluation.

Step-DAD introduces semi-amortized design: an offline policy is refined during
an experiment as observations arrive~\cite{hedman2025stepdad}. We retain that
infer--refine structure, using a posterior over parameters together with the
state carried by each posterior particle. A refinement is conditional on the
observed prefix and optimizes only future probes. This distinguishes online
adaptation from retraining a policy on the evaluator's latent parameter, which
would give the design method information unavailable in a real experiment.

\subsection{Objective-Based Uncertainty and Experimental Design}
MOCU measures uncertainty through the expected loss incurred when an operator
must be selected without knowing the true model~\cite{yoon2013mocu}.
Generalized MOCU extends the decision-oriented framework and its connection
to experimental design~\cite{boluki2018generalized}. This perspective explains
why reduction of parameter uncertainty is not, by itself, the same criterion
as reduction of decision loss. In the present application, the operator is a
continuous terminal injection and the finite loss penalizes its magnitude
and a model-dependent control shortfall. The specific surrogate loss and
posterior-coverage choice are declared here rather than attributed to a
universal monetary or reliability model.

Computational approximations are also important for objective-based design.
Chen et al.~\cite{chen2023objectiveoed} use neural message passing to accelerate
MOCU evaluation and experimental design for uncertain oscillator systems.
The current benchmark instead evaluates online trajectories and a numerical
terminal-control solver directly. A surrogate indexed only by the uncertain
parameters would not describe later non-reset stages unless it also represented
the relevant carried states and histories. No learned MOCU surrogate is used
in the comparisons specified here.

\section{Problem Formulation}
\label{sec:formulation}
\subsection{Power-System Model}
\label{sec:system-model}
\label{sec:ieee-system}
\subsubsection{Physical Network and Retained Machines}
The EIG study uses IEEE9 with nine electrical buses and three retained
dynamic machines, and IEEE14 with fourteen buses and five retained dynamic
machines. Their latent parameter dimensions are six and ten, respectively. Let $\mathcal N$ denote all physical
buses, $\mathcal M\subseteq\mathcal N$ the retained-machine buses, and
$m=|\mathcal M|$. The remaining buses are algebraic network nodes.
This distinction separates electrical topology from the dimension of the
dynamic model: an algebraic bus is not assigned an independent inertia or
frequency-response coefficient. The network determines how an injection at
a physical bus enters the retained-machine dynamics.

\subsubsection{Network Reduction and Injection Mapping}
Let $\mathcal A=\mathcal N\setminus\mathcal M$. For a connected lossless
network with unit voltage magnitudes, let
$\mathbf L$ be the weighted susceptance Laplacian, ordered as
$(\mathcal M,\mathcal A)$. If $\mathcal A\ne\varnothing$, assume
$\mathbf L_{\mathcal A\mathcal A}$ is invertible and define
\begin{align}
\mathbf L_{\mathrm{red}}&=\mathbf L_{\mathcal M\mathcal M}
-\mathbf L_{\mathcal M\mathcal A}\mathbf L_{\mathcal A\mathcal A}^{-1}
\mathbf L_{\mathcal A\mathcal M},\\
\mathbf A_{\mathrm{inj}}&=
\left[\mathbf I_m\; -\mathbf L_{\mathcal M\mathcal A}
\mathbf L_{\mathcal A\mathcal A}^{-1}\right],\\
B_{ij}&=-(L_{\mathrm{red}})_{ij}\quad(i\ne j),\qquad B_{ii}=0.
\label{eq:kron-map}
\end{align}
The physical-bus input map is $\mathbf a(b)=\mathbf A_{\mathrm{inj}}\mathbf e_b$
in the stated ordering. With no algebraic buses, use $\mathbf L_{\mathrm{red}}=\mathbf L$
and $\mathbf A_{\mathrm{inj}}=\mathbf I_m$. The linear Kron map combined with
sinusoidal coupling is a reduced-model approximation. Network data follow
the IEEE test cases in MATPOWER~\cite{zimmerman2011matpower}.


\subsection{Swing Dynamics and Latent Parameters}
\label{sec:dynamics}
\subsubsection{Inertia and Effective Frequency Response}
Motivated by the joint estimation setting of Peng et al.~\cite{peng2024probing},
we represent inertia and effective frequency response jointly in the following
reduced study model. Let $\Delta\omega(s)$ be angular-speed deviation
from an equilibrium with synchronous speed $\omega_{\mathrm s}=2\pi f_0$.
On a common power base $S_{\mathrm{base}}$, the incremental power balance is
\begin{equation}
M\frac{d\Delta\omega}{ds}
=\Delta P_m-\Delta P_e-D\Delta\omega,
\qquad M=\frac{2HS_n}{\omega_{\mathrm s}S_{\mathrm{base}}}.
\label{eq:inertia-balance}
\end{equation}
Here $H$ is the inertia constant, $S_n$ the rating associated with $H$,
$\Delta P_m$ and $\Delta P_e$ the input and electrical-output power changes,
and $D$ the damping coefficient. Power changes are in per unit.
The paper represents fast frequency support by a droop relation,
\begin{equation}
\Delta P_m=-K\Delta f=-\frac{K}{2\pi}\Delta\omega,
\qquad \Delta f=\frac{\Delta\omega}{2\pi},
\label{eq:frequency-support}
\end{equation}
where $K$ is expressed here in per-unit power per Hz. Substitution gives
\begin{equation}
M\frac{d\Delta\omega}{ds}
=-\Delta P_e-\left(\frac{K}{2\pi}+D\right)\Delta\omega.
\label{eq:aggregate_model}
\end{equation}
Thus $M$ governs the rate of frequency change, whereas $K$ governs the
restoring power supplied in response to a frequency deviation. We retain
known damping $D$ as an explicit component of the present reduced model;
Peng et al. neglect damping in their estimation model.

\subsubsection{Coupled Machine Dynamics}
Using the reduced network in Section~\ref{sec:ieee-system}, we extend the
aggregate response to the retained machines. For $i\in\mathcal M$, write $\omega_i$ for angular-speed
deviation and $\delta_i$ for rotor angle in the synchronous frame:
\begin{align}
\dot\delta_i(s)&=\omega_i(s),\\
M_i\dot\omega_i(s)&=-\Delta P_{e,i}(\boldsymbol\delta(s))
-\left(\frac{K_i}{2\pi}+D_i\right)\omega_i(s)+v_i(s).
\label{eq:swing-frequency}
\end{align}
The initial state is $(\boldsymbol\delta^0,\mathbf0)$, and the reduced
lossless-network power change is
\begin{equation}
\Delta P_{e,i}(\boldsymbol\delta)=\sum_{j\in\mathcal M}B_{ij}
\left[\sin(\delta_i-\delta_j)-\sin(\delta_i^0-\delta_j^0)\right],
\label{eq:network-power}
\end{equation}
with known coupling coefficients $B_{ij}$. A physical-bus injection is
mapped to $\mathbf v$ by the known vector $\mathbf a(b)$; the reduction is
given in Eq.~\eqref{eq:kron-map}. This machine-wise model is our extension
of the paper's aggregate response, with $K_i$ representing effective fast
frequency support at each retained machine bus.

\subsubsection{Latent Parameters and Prior Distribution}
The latent vector contains the inertia and response coefficients:
\begin{equation}
\boldsymbol\theta=(M_1,\ldots,M_m,K_1,\ldots,K_m)^\top\in\Theta,
\qquad\boldsymbol\theta\sim p_0,
\label{eq:theta-prior}
\end{equation}
where indices $1,\ldots,m$ follow the retained-machine ordering and
\begin{equation}
\Theta=\prod_{i=1}^{m}[M_i^-,M_i^+]
\times\prod_{i=1}^{m}[K_i^-,K_i^+],
\qquad M_i^->0,\quad K_i^-\ge0.
\label{eq:latent-space}
\end{equation}
The topology, coupling, damping, and equilibrium are known. Algebraic buses
enter the coupling and injection map but have no independent $M_i$ or $K_i$.
With bounded piecewise-continuous inputs, these assumptions define the
frequency trajectories used by the observation and terminal-control models.

\subsection{Probe Design and Observations}
\label{sec:design-problem}
\label{sec:probe-design}
\label{sec:problem}
\subsubsection{Probe Waveform and Non-Reset State Evolution}
The controllable probe variable is its duration, $\xi_t=\tau_t$.
Injection bus $b_{\mathrm p}$ and amplitude $A_{\mathrm p}$ are fixed.
A stage occupies a recording window of length $W$, independently of the
injection duration. With $s_t=(t-1)W$ and local time $\ell=s-s_t$, the
stage input is
\begin{equation}
\mathbf v_{\mathrm p,t}(\ell)
=\mathbf a(b_{\mathrm p})\frac{A_{\mathrm p}}{2}
\left[1-\cos\left(\frac{2\pi\ell}{\tau_t}\right)\right]
\mathbf1_{[0,\tau_t]}(\ell),\quad 0\le\ell\le W.
\label{eq:probe-input}
\end{equation}
Let $\mathbf x=(\boldsymbol\delta,\boldsymbol\omega)$ and let
$\Phi_\theta(\mathbf x,\tau;\ell)$ denote propagation of the swing equations
under this input. The state transition is
\begin{equation}
\mathbf x_t(\theta)=\Phi_\theta(\mathbf x_{t-1}(\theta),\tau_t;W),
\qquad \mathbf x_0=(\boldsymbol\delta^0,\mathbf0).
\label{eq:carry}
\end{equation}
Only the beginning of a new episode uses $\mathbf x_0$. Later stages use the
preceding terminal state. The latent $M_i,K_i$ do not change within an episode.
Consequently, $(\tau_1,\tau_2)$ and $(\tau_2,\tau_1)$ are different input
histories, even when their unordered duration sets coincide.

\subsubsection{Ordered Design Space}
Write $h_t=((\tau_1,\mathbf y_1),\ldots,(\tau_t,\mathbf y_t))$ for the
ordered available history. The design space at stage $t$ is
\begin{equation}
\Xi(h_{t-1})=\{\tau\in[\tau_{\min},\tau_{\max}]:
|\tau-\tau_k|\ge\varepsilon_\tau\text{ for every }k<t\},
\label{eq:design-space}
\end{equation}
where $0<\tau_{\min}<\tau_{\max}\le W$. This is a continuous set with
previously used neighborhoods excluded, not a catalog of bus--duration pairs.
For every horizon $T$ we restrict durations to increase:
$\tau_1<\cdots<\tau_T$. At stage $t$, the admissible interval is
$[a_t,\tau_{\max}-(T-t)\varepsilon_\tau]$, where
$a_1=\tau_{\min}$ and $a_t=\tau_{t-1}+\varepsilon_\tau$ for $t>1$.
This reserves room for subsequent probes. All methods select within this
interval before simulation; histories are never sorted after observing data.
Random samples uniformly within the current interval, not uniformly over
complete increasing sequences. This restriction applies at every horizon.


\subsubsection{Signed Endpoint-RoCoF Observation Model}
At measurement bus $b_{\mathrm y}$, the conference sensor reports
\begin{equation}
g_R(\theta,\mathbf x_{t-1}(\theta),\tau_t)
=\frac{\Delta f_{b_{\mathrm y}}(s_t+W)
-\Delta f_{b_{\mathrm y}}(s_t+W-\Delta r)}{\Delta r},
\label{eq:rocof-observation-functional}
\end{equation}
where $W=3.5$ s and $\Delta r=0.025$ s. Thus the frequency samples are
at 3.5 and 3.475 s relative to each probe start; no absolute value or
maximum is taken. Each candidate model starts from its own carried state.
The scalar observation model is
\begin{equation}
y_t=g_R(\theta,\mathbf x_{t-1}(\theta),\tau_t)+\epsilon_t,
\qquad\epsilon_t\sim\mathcal N(0,\sigma_R^2).
\label{eq:observation}
\end{equation}
Here $g_R$, $y_t$ and $\sigma_R$ have units Hz/s. Independent noise is
added to the clean summary, not to the physical state or individual frequency
samples. The explicit combination \texttt{endpoint\_rocof},
\texttt{N\_obs=0}, and $W=3.5$ selects this one-dimensional sensor.
The flag alone does not distinguish it from the separate maximum-absolute
RoCoF variant. Positive observation counts select a different sampled-frequency
model with noise in Hz. These observation protocols must not be pooled.

The scalar becomes available at window end. Computation and communication
latency are idealized as zero in simulated time; measured online runtime is
reported separately. A physical implementation with delay would need to
propagate the state during that delay.

\subsection{History-Conditioned Bayesian Inference}
The posterior is
\begin{equation}
p_t(\theta)=
\frac{p_{t-1}(\theta)
 p(\mathbf y_t\mid\theta,h_{t-1},\tau_t)}
{\int p_{t-1}(\vartheta)
 p(\mathbf y_t\mid\vartheta,h_{t-1},\tau_t)\,d\vartheta},
\label{eq:posterior}
\end{equation}
with the Gaussian likelihood from Eq.~\eqref{eq:observation}. Its mean must
be simulated from the candidate model's own carried state. An equilibrium
response indexed only by $(\theta,\tau_t)$ is insufficient after the first probe.

\subsection{Design Objectives}
\label{sec:objectives}
A useful experiment need not be the one that reduces parameter uncertainty
most strongly. For example, information that distinguishes models requiring
nearly identical controls can improve an information score without materially
changing the terminal action. Information that shifts posterior mass across
a control-requirement threshold can instead affect the decision directly.
This motivates evaluating three distinct utilities on the same probing
process rather than treating an information score as a proxy for control cost.

The non-reset setting adds a second mechanism: a probe changes both the
posterior and the physical state. Design objectives therefore depend on the
complete history, including the state from which control will act. We first
define parameter-oriented EIG, then specify the terminal-control task needed
to define MSC and finite-loss MOCU. The physical safety rate is a separate
evaluation quantity, not a substitute for either design objective.

\subsubsection{Expected Information Gain}
For a deterministic policy $\tau_t=\pi_\phi(h_{t-1})$, the information objective is
\begin{equation}
J_{\mathrm{EIG}}(\pi_\phi)=
\E_{\theta\sim p_0,\,H_T\mid\theta,\pi_\phi}
\left[\log\frac{p(\theta\mid H_T)}{p_0(\theta)}\right].
\label{eq:eig}
\end{equation}
The implementation uses the sequential prior-contrastive lower bound
(sPCE) of DAD~\cite{foster2021dad}, with history-conditioned likelihoods:
\begin{align}
\ell_t^{(n)}&=\sum_{k=1}^t
\log p(\mathbf y_k\mid\theta^{(n)},h_{k-1},\tau_k),\\
\widehat I_t&=\ell_t^{(0)}-
\log\left[\frac{1}{L+1}\sum_{n=0}^{L}\exp(\ell_t^{(n)})\right].
\label{eq:spce}
\end{align}
Here $\theta^{(0)}$ generates the observations and the other models are
independent prior draws. The reported expected bound depends on $L$ and
has ceiling $\log(L+1)$. The lower-bound statement concerns the expectation
of the Monte Carlo estimator, not a pointwise lower bound for every realized
trajectory. A finite-sample score can be negative. Historical bank-entropy
estimates represent a different statistic and are not combined with sPCE.

\subsubsection{Continuous Terminal Control and Physical Safety (Extension)}
The following control model is retained for MSC/MOCU extensions and is not
executed by the conference EIG campaign.
The final engineering decision is an injection magnitude
\begin{equation}
u_{\mathrm{ctrl}}\in\mathcal U=[u_{\min},u_{\max}],
\label{eq:control-space}
\end{equation}
which is continuous and distinct from the probe duration in
Eq.~\eqref{eq:design-space}. At global time $TW$, the probe input ends and
the declared contingency and supplementary step control begin. With local
control time $r$, the input over the control-assessment window $[0,H_c]$ is
\begin{equation}
\mathbf v_{\mathrm c}(r;u)=\mathbf p_{\mathrm c}+\mathbf a(b_{\mathrm c})u,
\qquad \mathbf x_{\mathrm c}(0;\theta,h_T)=\mathbf x_T(\theta).
\label{eq:control-input}
\end{equation}
The disturbance, injection location, window and units are fixed before design
optimization. Thus the connection from probing to action is explicit:
observations update the posterior over $M/K$, each candidate $M/K$ determines
a carried terminal state, and the controller evaluates safety from those states.

For the monitored dynamic buses $\mathcal M$, define
\begin{align}
f_{\min}(\theta,h,u)&=\min_{i\in\mathcal M,\,0\le r\le H_c}
\Delta f_i(r;\theta,\mathbf x_T(\theta),u),\\
R_{\max}(\theta,h,u)&=\max_{i\in\mathcal M,\,0\le r\le H_c}
\left|\frac{d\Delta f_i(r;\theta,\mathbf x_T(\theta),u)}{dr}\right|,\\
S(\theta,h,u)&=\mathbf1\{f_{\min}(\theta,h,u)\ge\underline f,
\ R_{\max}(\theta,h,u)\le\overline R\}.
\label{eq:safety}
\end{align}
The expressions in Eq.~\eqref{eq:safety} define the continuous-time target.
The implemented control-safety calculation uses the RK4 trajectory at step
$\Delta s=1/640$~s, with $J_c=H_c/\Delta s$, and computes
\begin{align}
\widehat f_{\min}&=\min_{i\in\mathcal M,\,j=0,\ldots,J_c}
\Delta f_i(j\Delta s),\\
\widehat R_{\max}&=\max_{i\in\mathcal M,\,j=1,\ldots,J_c}
\left|\frac{\Delta f_i(j\Delta s)-\Delta f_i((j-1)\Delta s)}{\Delta s}\right|.
\label{eq:discrete-control-safety}
\end{align}
The initial frequency is included in the nadir and in the first RoCoF
difference. These safety metrics are evaluated at every integration step;
the 40~Hz differencing interval specifies the probe observation only.
The numerical safety indicator compares these computed extrema with the
same declared thresholds. It approximates the continuous-time target and
does not certify unsampled extrema or general transient stability throughout
the preceding probes.
The model-specific minimum safe control is
\begin{equation}
U(\theta,h)=\inf\{u\in\mathcal U:S(\theta,h,u)=1\},
\label{eq:oracle}
\end{equation}
with $U=+\infty$ when the admissible safe set is empty.
The current scalar solver assumes that a nonempty safe set is
$[U(\theta,h),u_{\max}]$. It checks sampled monotonicity and rejects a detected
violation; those numerical checks do not prove the assumption globally.

\subsubsection{Minimum Safe Control as a Design Objective}
For a specified posterior coverage $q\in(0,1)$, select
\begin{equation}
u_{\mathrm{MSC}}(h;q)=\inf\left\{u\in\mathcal U:
\int S(\theta,h,u)p(\theta\mid h)\,d\theta\ge q\right\}.
\label{eq:msc-controller}
\end{equation}
If this set is empty, the decision is infeasible. Under the scalar safety
assumption, Eq.~\eqref{eq:msc-controller} is the posterior $q$ quantile of
$U(\theta,h)$. Coverage expresses a conditional model-based requirement;
it is not an observed safety rate or a value established by an engineering standard.

We use \emph{MSC} as a descriptive abbreviation for minimum safe control,
not as a claim of standardized terminology. Its BOED target is
\begin{equation}
\min_\pi J_{\mathrm{MSC}}(\pi),\qquad
J_{\mathrm{MSC}}(\pi)=\E_{H_T\mid\pi}[u_{\mathrm{MSC}}(H_T;q)].
\label{eq:msc-objective}
\end{equation}
The physical requirement $U(\theta^\star,h)$ evaluated with true parameters
is the oracle; $u_{\mathrm{MSC}}(h;q)$ is the implementable decision using
noisy observations. They are not interchangeable. Lower expected MSC values
smaller selected control at the declared posterior coverage, without an
under-control penalty in the design score.

\subsubsection{Finite-Loss Control-Oriented MOCU}
% BEGIN FINITE-LOSS MOCU
MOCU compares the loss of a decision made under uncertainty with that of a
model-informed decision~\cite{yoon2013mocu,boluki2018generalized}.
Our particular control-sizing loss, for finite feasible $U(\theta,h)$, is
\begin{equation}
C_q(u,\theta;h)=u+\frac{[U(\theta,h)-u]_+}{1-q}.
\label{eq:operational_cost}
\end{equation}
This is a chosen surplus/shortfall surrogate in control-magnitude units.
It is not an engineering-calibrated cost of violating the frequency or RoCoF
limits. The cited MOCU framework supplies the decision-theoretic construction,
not this particular loss or an engineering value of $q$.

For fixed $h$, the model-informed optimum is $u=U(\theta,h)$ with loss
$U(\theta,h)$. A posterior Bayes action satisfies
\begin{equation}
u_q(h)\in\arg\min_{u\in\mathcal U}
\E_{\theta\mid h}[C_q(u,\theta;h)].
\label{eq:bayes-action}
\end{equation}
Writing $F_h$ for the posterior CDF of $U(\theta,h)$, its optimality condition is
\begin{equation}
F_h(u_q^-)\le q\le F_h(u_q).
\label{eq:quantile_optimality}
\end{equation}
Choosing the smallest such action gives the same controller as
Eq.~\eqref{eq:msc-controller} under the stated threshold assumption. The
\emph{design objectives} nevertheless differ: the posterior MOCU is
\begin{equation}
\mathcal M(h;q)=\E_{\theta\mid h}
\left[u_q(h)+\frac{[U(\theta,h)-u_q(h)]_+}{1-q}-U(\theta,h)\right],
\label{eq:mocu}
\end{equation}
and MOCU-based design minimizes
\begin{equation}
J_{\mathrm{MOCU}}(\pi)=\E_{H_T\mid\pi}[\mathcal M(H_T;q)].
\label{eq:mocu-objective}
\end{equation}
MSC scores selected control; MOCU scores expected excess loss relative to
the model-specific requirement. MOCU is nonnegative under its assumptions
and vanishes when the required control is known exactly. A lower MOCU need
not imply a smaller selected control or a larger empirical safety rate.
Infeasible requirements make this finite-loss construction unsupported and
must not be hidden by clipping or resampling.

In the non-reset application, probing changes both the posterior and the
physical state entering control. Thus $U(\theta,h)$ itself depends on the
applied probe history. A fixed-application value-of-information monotonicity
argument cannot be transferred without qualification: increasing the horizon
need not reduce either MSC or MOCU, and any improvement is not automatically
attributable to information alone.
% END FINITE-LOSS MOCU

\section{Methodology}
\label{sec:methodology}
\subsection{Common Information State and Feasible Actions}
Let $h_{t-1}$ contain the ordered durations and observations available before
probe $t$. For a fixed horizon, the policy input is
\begin{equation}
z_t=\left[\frac{t-1}{T},
 \left(m_k\frac{\tau_k-\tau_{\min}}{\tau_{\max}-\tau_{\min}},
 m_k\frac{\mathbf y_k}{s_y},m_k\right)_{k=1}^{T}\right],
\label{eq:history-features}
\end{equation}
where $m_k=1$ for $k<t$ and is zero otherwise; unobserved slots are zero.
The observation scale is $s_y=0.05$ in the units of the chosen observation
model. Thus the default scalar observation gives $1+3T$ input coordinates.
The stage index and masks distinguish an unobserved entry from a measured
zero. No latent parameter, true terminal state, or future observation is an
input to the design policy.

Let $a_t$ and $b_t=\tau_{\max}-(T-t)\varepsilon_\tau$ be the endpoints of
the increasing feasible interval defined in Section~\ref{sec:design-problem}.
All methods first select a unit coordinate $v_t\in[0,1]$ and then apply
\begin{equation}
\tau_t=a_t+(b_t-a_t)v_t.
\label{eq:common-action-map}
\end{equation}
This enforces duration order and reserves room for every remaining stage
before simulation. The implementation adds a small numerical margin to the
minimum separation to avoid roundoff violations. It neither sorts completed
actions nor snaps them to a duration grid. Random and Myopic know the remaining
feasible interval even though they do not optimize a full-horizon objective.

\begin{table}[t]
\caption{Adaptation and planning in the implemented design strategies.
Adaptive means using observations from the current experiment to choose
later probes. Myopic means optimizing only the next probe's expected utility.
N/A denotes no utility optimization. MoE-sBOED is optional and excluded from the primary six-method campaign.}
\centering
\begin{tabular}{lcc}
\toprule
Method & Adaptive? & Myopic?\\
\midrule
Random & No & N/A\\
Fixed & No & No\\
Myopic & Yes & Yes\\
DAD & Yes & No\\
RL-sBOED & Yes & No\\
Step-DAD & Yes & No\\
MoE-sBOED (optional) & Yes & No\\
\bottomrule
\end{tabular}
\label{tab:method-capabilities}
\end{table}

Fixed learns a shared probing sequence offline and does not adapt it to
observations during evaluation. DAD and RL-sBOED use trained policies to
respond to observations without updating their weights during an experiment.
Step-DAD additionally updates its policy after observations using simulated
future trajectories conditioned on the observed history.

For each objective, define the terminal score
\begin{equation}
G_T=\begin{cases}
\widehat I_T, & \text{EIG},\\
-u_{\mathrm{MSC}}(h_T;q), & \text{MSC},\\
-\widehat{\mathcal M}(h_T;q), & \text{MOCU}.
\end{cases}
\qquad J(\phi)=\E[G_T].
\label{eq:common-terminal-utility}
\end{equation}
The expectation includes parameter and measurement-noise draws and, when
used, stochastic policy actions. Each objective is trained separately; an EIG
checkpoint is not used as an MSC/MOCU policy.
All methods share the physical model and observation likelihood; MSC and
MOCU additionally use the same declared terminal-control assumptions.
An EIG-only evaluation reports information and does not itself produce a
terminal-control safety assessment.

\subsection{DAD: Offline Full-Horizon Policy Optimization}
The deterministic policy is
$v_t=\operatorname{sigmoid}(f_\phi(z_t)+c_t)$, where $c_t$ is a trainable
stage bias. The implemented network has two hidden layers of width 64 with
tanh activations and a scalar output. This ordered representation retains the
policy-based construction of DAD~\cite{foster2021dad} while respecting the
non-reset physical model. Training maximizes the Monte Carlo estimate of
Eq.~\eqref{eq:common-terminal-utility}; evaluation uses the selected policy
without changing its weights during an episode.

For EIG, parameters and standard Gaussian observation noises are sampled,
and complete trajectories are generated with the same latent model carried
across stages. The sPCE score is computed using the truth-generating model and
$L$ independent prior contrast models as in Eq.~\eqref{eq:spce}. Every contrast
model follows the actual durations selected from the generated history, but
propagates its own physical state. Gradients pass through the policy, feasible
action map, observations and carried states over the full horizon.

In a joint EIG run, DAD compares a random initialization against a policy
initialized from that run's newly optimized Fixed sequence. For the latter,
the final network layer is set to zero and the sequence logits populate the
stage biases after clipping to $[-4,4]$. This creates a trainable interior
initial policy; it need not reproduce a boundary-valued Fixed sequence exactly.
The choice uses a fixed validation simulation set, never evaluation outcomes.
Fixed training is a prerequisite for this initialization, so its cost must
be disclosed when reporting the total cost of that DAD training procedure.
This describes fresh training; evaluating a previously trained matching
checkpoint does not repeat initialization or training.

\subsection{RL-sBOED: Incremental Rewards with Full-Horizon Return}
Let $G_t$ denote the objective score assessed after prefix $h_t$, with $G_0=0$.
Transitions receive
\begin{equation}
r_t=G_t-G_{t-1},\qquad \gamma=1,
\qquad \sum_{t=1}^{T}r_t=G_T.
\label{eq:rl-rewards}
\end{equation}
For EIG, this is the increment of sPCE, not an exact information gain available
without simulation. For MSC/MOCU, prefix scores use hypothetical terminal
controls at that prefix; these controls are reward calculations, not additional
physical interventions. The undiscounted unregularized return is therefore
the same terminal criterion used by the deterministic methods. A myopic
optimizer would maximize only the next increment, whereas the RL critic
includes future returns.

The implemented RL-sBOED adaptation uses a tanh-Gaussian actor with two
128-unit ReLU hidden layers, a replay buffer and a REDQ-style critic
ensemble~\cite{blau2022rlsboed}. If $a_t=\tanh(\zeta_t)\in[-1,1]$ is the
normalized stochastic action, the feasible-map coordinate is $v_t=(a_t+1)/2$.
The log density includes the tanh Jacobian and is measured in $a_t$ coordinates,
while critics receive $v_t$. These conventions matter for entropy tuning.

For a replay transition $(z,v,r,d,z')$, where $d$ marks termination, the
critic target and actor loss are
\begin{align}
y_Q&=r+(1-d)\left[
 \min_{j\in\mathcal S}\overline Q_j(z',v')-\alpha\log\pi_\phi(a'\mid z')\right],\\
\mathcal L_\pi&=\E\left[\alpha\log\pi_\phi(a\mid z)
 -\frac{1}{N_Q}\sum_{j=1}^{N_Q}Q_j(z,v)\right].
\label{eq:redq-objectives}
\end{align}
Here $\mathcal S$ is a uniformly chosen pair of target critics; each critic
minimizes squared error to $y_Q$. The current settings use ten critics,
a replay capacity of $10^5$ transitions, replay minibatches of 256 and five
updates per newly simulated trajectory batch. Target parameters move toward
the online parameters with coefficient 0.005. Actor, critic and temperature
optimizers use Adam with learning rate $3\times10^{-4}$.

The initial temperature is 0.01, and its target entropy is $-1$ for the
one-dimensional normalized action. The temperature objective is
$\mathcal L_\alpha=-\E[\log\alpha\,
(\log\pi_\phi(a\mid z)-1)]$, treating the sampled log density as constant
for this update. This regularizer changes the training objective; entropy
bonuses are excluded from reported EIG/MSC/MOCU scores. Evaluation transforms
the Gaussian mean deterministically through the same action map. Recorded
entropy, temperature, critic error and target statistics are diagnostics,
not proofs of policy convergence.

\subsection{Step-DAD: Posterior-Conditional Online Refinement}
Step-DAD begins each episode with the DAD checkpoint trained in the same run.
Its first probe uses that policy unchanged. Once a prefix is observed, each
planning particle is propagated under the realized durations and reweighted
by the likelihood. A private policy copy is then refined for the remaining
horizon, following the semi-amortized structure of
Hedman et al.~\cite{hedman2025stepdad}.

For EIG refinement, a fantasy truth and its contrasts are sampled from the
current discrete posterior, with their corresponding carried states. Future
observations are simulated from the fantasy truth, the observed prefix is held
fixed, and the conditional future sPCE objective replaces a restart from the
original prior. The actual evaluation truth is never supplied to this procedure.
MSC/MOCU refinements instead use their terminal control utility from these
posterior-conditional simulations.

Each refinement call creates a fresh optimizer and performs the declared
number of updates. Candidate policies are compared on a separate fixed set of
conditional validation simulations; the starting policy is retained if no
candidate improves that estimate. Policy weights selected at one stage seed
the next refinement within the episode, while a new episode starts again from
the original DAD checkpoint. Thus the method carries adapted weights, not an
optimizer state or observations from other evaluation systems. Improvement
on conditional validation does not guarantee improvement on the eventual
realized trajectory.

\subsection{MoE-sBOED (Optional Extension)}
\label{sec:optional-moe}
\textbf{Optional:} the current power-grid implementation is a deterministic,
history-conditioned soft mixture, separate from the six-method conference
campaign. Each of four experts receives the ordered features $z_t$ in
Eq.~\eqref{eq:history-features}. Each expert has two 64-unit tanh hidden
layers and a scalar latent-action output, with its own trainable stage bias.
A router with one 64-unit tanh hidden layer produces four logits. With
expert parameters $\psi_e$, router parameters $\eta$, and stage biases
$c_{e,t}$, the action is
\begin{align}
w_e(z_t)&=\frac{\exp r_{\eta,e}(z_t)}
{\sum_{j=1}^{4}\exp r_{\eta,j}(z_t)},\\
u_t^{\mathrm{lat}}&=\sum_{e=1}^{4}w_e(z_t)
\bigl[f_{\psi_e}(z_t)+c_{e,t}\bigr],\\
v_t&=\operatorname{sigmoid}(u_t^{\mathrm{lat}}),\qquad
\tau_t=a_t+(b_t-a_t)v_t.
\label{eq:optional-moe-action}
\end{align}
The mixture acts on latent actions before the sigmoid and common feasible
map. All four experts execute for each input: this is dense soft routing,
not top-$k$ sparse execution. It is distinct from the separate SIR sparse-MoE
implementation. IEEE9 and IEEE14 share this architecture but train separate
weights under their own physical models.

Experts and router are jointly trained offline to maximize full-horizon sPCE
with the same numerical pathwise simulator derivatives used for DAD.
The router initially assigns equal weights. A declared Fixed-sequence
initialization may set the expert outputs to a common interior sequence;
this is initialization rather than a DAD teacher or imitation objective.
Its source and computational cost must be reported. During evaluation,
weights remain fixed: the mixture responds to observations without the
online gradient refinement used by Step-DAD. No planner receives the held-out
true parameters or true carried state.

Different expert outputs or approximately balanced router weights alone do
not establish useful specialization. An optional evaluation should report
router variation across histories, expert-output differences, online runtime,
and comparisons with uniform routing and a parameter-matched single network,
using paired evaluation cases and independent training seeds. Such ablations
are proposed evidence, not completed results. No superiority over DAD is
claimed, and MoE-sBOED scores must remain separate from the primary tables.

\subsection{Fixed, Myopic and Random Baselines}
Fixed uses trainable stage logits $c_1,\ldots,c_T$ and
$v_t=\operatorname{sigmoid}(c_t)$, independent of observations. The sequential
map in Eq.~\eqref{eq:common-action-map} makes this a deterministic increasing
sequence. Initialization candidates cover short, central, long and spread
interior durations. Each desired sequence is converted to feasible-interval
logits, and validation selects the starting candidate. Initial logits such as
$\pm20$ are avoided because sigmoid saturation can suppress useful gradients.
Fixed then receives full-horizon optimization and validation-based checkpoint
selection; initialization selection alone is not reported as successful
optimization.

Myopic selects a duration by maximizing expected one-step utility under the
current particle belief. Its EIG criterion is estimated by
\begin{equation}
\widehat V_t(\tau)=\frac{1}{F}\sum_{f=1}^{F}
\log\frac{p(y^{(f)}\mid\theta^{(f)},h_{t-1},\tau)}
{\sum_n w_n p(y^{(f)}\mid\theta^{(n)},h_{t-1},\tau)},
\label{eq:myopic-utility}
\end{equation}
where fantasy models follow the current posterior and each model uses its
own carried state. For MSC/MOCU, fantasies instead score the next posterior's
control objective. Search uses continuous random proposals followed by elite
mean/variance refinement in the unit interval. Common random draws are used
across proposals within a decision, and the best observed proposal is retained.
This finite search is not a global one-step optimum. It does not propagate a
lookahead tree to the end of the remaining horizon.

Random draws $v_t$ uniformly from $[0,1]$ at each stage. It is uniform over the
current feasible duration interval, not over all complete increasing sequences.
Neither Random nor Myopic has an offline policy-training cost, but Myopic
incurs online posterior and candidate-evaluation costs. All baseline scores
are evaluated with the same final estimator used for the learned methods.

\subsection{Online Posterior and Continuous Control Solver}
For posterior particles $\theta^{(n)}$, all states are propagated under the
actual ordered durations. Weights use
\begin{equation}
w_n^{(t)}\propto w_n^{(t-1)}
p(\mathbf y_t\mid\theta^{(n)},h_{t-1},\tau_t),\qquad \sum_nw_n^{(t)}=1.
\label{eq:particle-update}
\end{equation}
The particle support is fixed within an episode; weights and states are updated without particle rejuvenation. Effective sample size is monitored because small support can become unrepresentative after several probes. Increasing the support is a sensitivity analysis, not a correction obtained by inserting the evaluator truth.
The controller's support is sampled independently of the system generating
an evaluation trajectory; that generating model is not deliberately inserted
into this posterior. Its role in the sPCE denominator is specific to
Eq.~\eqref{eq:spce} and must not be transferred to control inference.

For each carried state, a numerical scan of the discretized safety indicator
brackets the first detected
unsafe-to-safe transition, followed by bisection to approximate
$U_n=U(\theta^{(n)},h)$. The scan points are not admissible-action restrictions:
final controls take continuous values between them. A detected non-monotone
safe set stops the scalar solver. With normalized weights,
\begin{align}
\widehat F_h(v)&=\sum_n w_n\mathbf1\{\widehat U_n\le v\},\\
\widehat u_q(h)&=\inf\{v\in\mathcal U:\widehat F_h(v)\ge q\},\\
\widehat{\mathcal M}(h;q)&=\sum_nw_n
\left[\widehat u_q+\frac{(\widehat U_n-\widehat u_q)_+}{1-q}-\widehat U_n\right].
\label{eq:particle-objectives}
\end{align}
The selected action is checked again by direct safety simulation across the
posterior support when computing terminal decisions and full-horizon training
scores. Myopic control fantasies use the weighted numerical requirements
without an additional direct safety simulation for each fantasy posterior. For MSC, an infeasible particle keeps its original
posterior mass with $\widehat U_n=+\infty$; a finite decision exists only if
coverage can still be met. Finite-loss MOCU instead requires finite feasible
requirements. No failed system is discarded to improve reported safety.

For MSC/MOCU, the full evaluation records ordered actions, noisy observations, posterior
weights and states, selected continuous control, true-system safety and a
separately computed continuous numerical oracle at the actual terminal state.
An equilibrium control bank indexed only by $M/K$ cannot replace this calculation.

\subsection{Pathwise Simulator Gradients}
The conference endpoint sensor differentiates the signed difference of the
two fixed terminal frequency samples in
Eq.~\eqref{eq:rocof-observation-functional}. It has no peak-selection operation.
The branch-consistency discussion below applies only to the separate
maximum-absolute-RoCoF observation variant.
Write one simulator stage as
$(\mathbf x_t,g_t)=\mathcal F(\theta,\mathbf x_{t-1},\tau_t)$.
For deterministic EIG policies, differentiation must include both the action
path and the carried-state path. In particular,
\begin{equation}
\frac{d\mathbf x_t}{d\phi}=
\frac{\partial\mathcal F_x}{\partial\mathbf x_{t-1}}
\frac{d\mathbf x_{t-1}}{d\phi}
+\frac{\partial\mathcal F_x}{\partial\tau_t}
\frac{d\tau_t}{d\phi}.
\label{eq:state-chain-gradient}
\end{equation}
The action derivative also includes dependence on earlier observations and
the moving lower endpoint of the feasible interval. Differentiating only the
current duration while treating the incoming state as independent would omit
how earlier probes change later information.

Simulator Jacobians are approximated numerically, using central differences
with state perturbation $10^{-5}\max(1,|x_j|)$ and duration perturbation
$10^{-5}$~s, clipped at the duration bounds. Let $r_k$ denote a sampled signed
frequency difference divided by its sampling interval. At a trajectory with a
unique maximizing index $k^*=\arg\max_k|r_k|$ and nonzero $r_{k^*}$,
\begin{equation}
D\max_k|r_k|=\operatorname{sign}(r_{k^*})D r_{k^*}.
\label{eq:active-rocof-gradient}
\end{equation}
The implementation records this index and sign on the unperturbed trajectory
and holds them fixed while forming observation Jacobians. The forward
observation remains the maximum over all samples. Independently reselecting
a maximum for every coordinate perturbation can mix different branches and
produce a Jacobian inconsistent with the original trajectory. Peak ties remain
nondifferentiable; a fixed-branch numerical derivative does not remove that
limitation or establish global smoothness.

Verification compares full sPCE directional derivatives with independent
end-to-end finite differences for both Fixed and adaptive policies at
$T=3,4,5$. Conditional-prefix checks exercise the gradients used by Step-DAD.
Forward calculations must agree with the explicitly selected observation model.
These checks concern the derivative implementation; training histories and
sensitivity to numerical scales remain necessary for assessing optimization.

\subsection{Numerical Optimization of Control Objectives}
A continuous control variable does not make a finite-particle empirical
quantile smooth in posterior weights. MSC/MOCU training therefore retains
deterministic design policies but uses paired, antithetic parameter
perturbations, following the black-box optimization principle of
Salimans et al.~\cite{salimans2017es}. For independent
$z_k\sim\mathcal N(0,I)$ and scale $d>0$, the estimator is
\begin{equation}
\widehat{\nabla J_d}(\phi)=\frac1K\sum_{k=1}^K
\frac{\widehat J(\phi+d z_k)-\widehat J(\phi-d z_k)}{2d}\,z_k.
\label{eq:numerical-gradient}
\end{equation}
Each positive/negative pair uses the same model and noise draws. Every utility
evaluation still calls the original numerically resolved control objective;
no differentiable replacement of the safety indicator is substituted.
The target is the Gaussian-smoothed objective
$J_d(\phi)=\E_z[J(\phi+dz)]$, with finite-smoothing bias relative to $J$.
Validation and evaluation use the unperturbed policy. Scale and direction-count
sensitivity are required before interpreting a performance difference.
This is a stated numerical adaptation of the deterministic methods; it is
neither the original pathwise DAD estimator nor a stochastic-action policy
gradient. RL-sBOED avoids simulator and terminal-controller derivatives by
using observed simulated rewards in its actor--critic updates.

\subsection{Training, Checkpoint Selection and Reproducibility}
Offline training uses independent simulated batches and a fixed validation
set, with checkpoints selected by the relevant unregularized validation
utility. EIG DAD/Fixed use Adam with learning rate $10^{-3}$ and gradient-norm
clipping at 5; Step-DAD uses the same learning rate for its conditional updates.
For control objectives, numerical parameter gradients use the same clipping
threshold. Equal update counts do not imply equal simulation counts or costs
across these optimizers.

Evaluation resets the random-number generator for each method so that the
latent systems, prior contrast models and measurement-noise draws are paired.
Method-specific actions still induce different observations and terminal
states. Planning simulations use a separate random stream and particle support.
The evaluation truth can appear in the sPCE scoring denominator but cannot
enter a planner or terminal controller by privileged access. Each independent
training seed defines a fresh policy fit. Additional evaluation seeds may
explicitly reuse completed checkpoints with the same training seed, horizon,
objective, physical and observation settings, training budgets and scientific
implementation. Fresh training remains the default. Such evaluations retain
the original checkpoint provenance and do not count as independent training
runs. Step-DAD always initializes from the corresponding DAD checkpoint.

For each trained method, record initialization scores, the selected update,
full validation history, gain over initialization and recent validation change.
Assess whether improvements persist near the budget limit, whether the selected
checkpoint is stable, and whether the result is sensitive to optimization and
Monte Carlo budgets. No fixed recent-change threshold by itself proves
convergence. Report offline training, online decision/refinement and
objective-specific evaluation costs separately, together with the execution
hardware. Numerical correctness, practical convergence and evidence across
independent seeds are distinct requirements.

\section{Experimental Setup}
\label{sec:experimental-protocol}
\subsection{Conference Probe Settings and Control Extensions}
IEEE9 has three retained machines and six uncertain $M_i/K_i$ coordinates;
IEEE14 has five machines and ten coordinates. Both EIG campaigns use fixed
physical probe and measurement bus 1, amplitude $0.05$ pu, continuous
durations in $[0.2,3]$ s, increasing order with minimum separation $0.01$ s,
and $W=3.5$ s. Horizons $T=3,4,5$ therefore span 10.5, 14 and 17.5 s of
simulated probing. Each window yields the signed endpoint-RoCoF observation
in Eq.~\eqref{eq:rocof-observation-functional}, with independent noise
$\sigma_R=0.005$ Hz/s. Classical fixed-step RK4 uses step $1/640$ s;
frequency differencing uses 40 Hz samples. Every machine state carries
between windows, and durations remain continuous.

The EIG campaign does not execute a terminal-control solver or report
uncomputed safety metrics. Retained MSC/MOCU examples use a continuous
control on $[0,0.5]$ pu, a $-0.45$ pu contingency at bus 1, a 10 s control
window, coverage $q=0.9$, and specified nadir/RoCoF limits. These are
separate study assumptions, not operating standards or settings inferred
from the EIG results.

\subsection{IEEE9 Prior Specification}
\begin{table}[t]
\caption{IEEE9 independent uniform prior bounds and known damping. Values are rounded here; resolved YAML values are retained with each run. Units follow Eq.~\eqref{eq:swing-frequency}.}
\centering
\begin{tabular}{lrrr}
\toprule Parameter & Machine 1 & Machine 2 & Machine 3\\
\midrule
$M_i^-$ & 0.0877899 & 0.0237671 & 0.0111780\\
$M_i^+$ & 0.1630383 & 0.0441390 & 0.0207591\\
$K_i^-$ & 0.029375 & 0.013125 & 0.007500\\
$K_i^+$ & 0.293750 & 0.131250 & 0.075000\\
$D_i$ & 0.0125414 & 0.0033953 & 0.0015969\\
\bottomrule
\end{tabular}
\end{table}
The prior is uniform on the product set in Eq.~\eqref{eq:latent-space};
its coordinate widths are study choices, not identification results.

\subsection{Operating Point and Numerical Settings}
The benchmark begins with zero relative rotor angles and zero speed deviations,
with nominal incremental power balance consistent with that equilibrium.
This operating point is a reduced perturbation-model assumption, not a
validated loaded AC dynamic operating point.

Each run records its resolved priors, coupling, damping, simulator settings
and seeds. Coverage and physical thresholds are specified before final testing.
The online benchmark does not use equilibrium observation or control banks.

\subsection{Training and Planning Budgets}
The specified full-budget EIG comparison uses the explicitly submitted
configuration rather than unmodified template defaults. In particular, the
batch size is overridden from the default 64 to 32. Each offline-trained
method receives 2,000 trajectory-batch updates,
batch size 32, 128 contrast models, and 128 validation systems evaluated every
100 updates and at initialization. Evaluation uses 512 physical systems per evaluation
seed, with three evaluation seeds per trained policy (1,536 episodes) and
128 independent planning particles. Myopic uses 16 candidate
proposals per round, three rounds and 32 fantasies per proposal. Step-DAD uses
four conditional refinement updates at each stage after the first,
with batch size 32 and validation after each update. This gives 8, 12 and
16 total gradient updates per episode at $T=3,4,5$, respectively. The earlier
16-update Step-DAD campaign is superseded and excluded from the conference comparison. These are
computational settings, not established convergence thresholds. The
conference contrast count is $L=128$, so its sPCE ceiling is
$\log(129)\simeq4.8598$ nats. Separate MoE development runs may use
$L=1024$ with ceiling $\log(1025)\simeq6.9324$ nats; those scores must not
be pooled with this campaign. The EIG-only
runs do not call the MSC/MOCU terminal-control solver, even though the shared
configuration retains its settings. Control-objective extensions use numerical
parameter perturbations (default scale 0.01 and four directions per update);
those budgets and any changes require separate reporting and sensitivity
analysis.

\subsection{Physical Safety Rate and Held-Out Evaluation}
The policy safety probability is
\begin{equation}
\mathrm{SR}(\pi;q)=\E_{\theta\sim p_0,\,H_T\mid\theta,\pi}
[S(\theta,H_T,u_q(H_T))].
\label{eq:safety-rate}
\end{equation}
With an exact posterior and the same correctly specified prior-predictive
model, conditional coverage in Eq.~\eqref{eq:msc-controller} implies
$\mathrm{SR}(\pi;q)\ge q$ by iterated expectation. A finite-particle
approximation, model mismatch or a finite test sample does not provide that
empirical guarantee.
For $N$ independent test systems and $R_i$ observation-noise replicas per system,
\begin{equation}
\widehat{\mathrm{SR}}=\frac1N\sum_{i=1}^N\frac1{R_i}
\sum_{r=1}^{R_i}S(\theta_i^\star,h_{i,r},u_q(h_{i,r})).
\label{eq:empirical-safety}
\end{equation}
Controls and posterior objectives use the same within-system averaging.
The held-out oracle is recomputed using $\theta_i^\star$ and that method's
actual terminal state. Different probe histories can therefore have different
oracle values even for the same true $M/K$. True parameters are used only
to generate/evaluate physical outcomes, never as controller or policy inputs.

\subsection{Verification and Comparison Scope}
A workability test uses $T=3$, training seed 101, evaluation seed 1001,
two training updates, training batch size four, four validation systems, two test systems,
eight posterior/planning particles and all six methods. It checks execution,
state carry, objective consistency and output completeness; its scores are
not performance evidence. Independent CPU/GPU trajectory checks, continuous
root-solver checks, posterior-truth exclusion and RL reward telescoping are
separate implementation checks.

The submitted IEEE9 and IEEE14 EIG protocol uses $T=3,4,5$, training
seeds 101/202/303 and evaluation seeds 1001/1002/1003. These are three
independent trainings, not nine. Each trained policy is evaluated on
512 systems per evaluation seed. Random and Myopic have no training-seed
dependence under the fixed evaluation streams; their identical evaluations
may be computed once per horizon and referenced across training seeds.
Fixed remains separately trained for each training seed. Step-DAD reuses
the corresponding DAD checkpoint, with its explicit four-update online
budget recorded in evaluation provenance.

Evaluation seed 1001 has informed development and is not an untouched
confirmation set. Conclusions must distinguish exploratory evidence from
confirmation after methodological choices are fixed. Contrast and particle
sensitivity, numerical checks and checkpoint convergence are separate from
simply completing the scheduled budget. No ranking is assumed in advance.

Each metric gets a separate table with methods as rows and horizons as columns.
Average the three evaluation-seed means within each training seed, then
report the mean and sample standard deviation across the three training seeds;
per-system Bernoulli dispersion is not the uncertainty of a reported safety
rate. A single evaluation seed has no across-seed standard deviation.
Keep physical systems and their repeated noise outcomes together in paired
uncertainty calculations. Fixed training/calibration copies may be deduplicated
only if the resulting policies and evaluations are actually identical.
Report posterior MSC or MOCU as the primary objective-specific score, with
control, physical safety and the history-dependent oracle as appropriate.
Realized regret and posterior MOCU are treated as distinct statistics. Timing must distinguish
offline optimization, online probe choice/refinement, terminal posterior/control
computation and evaluation-only oracle work.

\subsection{Evaluation Metrics and Comparison Integrity}
For EIG, the primary statistic is the mean terminal sPCE over the $N$ paired
physical systems for a given training/evaluation seed pair,
\begin{equation}
\overline I_{\mathsf A,T}=\frac{1}{N}\sum_{i=1}^{N}\widehat I_{\mathsf A,T}^{(i)},
\label{eq:reported-spce}
\end{equation}
where $\mathsf A$ denotes the design method. The contrast count and the ceiling
$\log(L+1)$ accompany the score. Dispersion across physical systems measures
heterogeneity within that run; it is not a standard deviation across
independent training runs. With a single seed pair, the latter is unavailable.
When comparing two methods, differences are formed for the same physical
systems before any paired uncertainty calculation. A bootstrap over systems,
if reported, is conditional on the trained policies and does not measure
variation due to retraining.

For MSC and MOCU, report the corresponding terminal objective separately from
selected control magnitude, safe-system counts and history-dependent oracle
values. Nats, control magnitudes and loss values are not combined into a joint
numerical ranking. Likewise, the scalar probing observation must not
be confused with the monitored safety quantities over the terminal-control
window. An EIG-only experiment cannot supply uncomputed control-safety metrics.

A run enters the comparison only after the requested methods and evaluation
systems have completed and its per-system records agree with its summary.
Acceptance checks require finite objective values, duration bounds and
ordering, common physical and observation settings, and implementation
provenance. Failure records are retained and missing cells are reported;
failed systems are not removed or replaced by results from another protocol. Training and
validation budgets, selected checkpoints and execution hardware accompany
each accepted result. This establishes traceability; it does not certify
convergence or compensate for insufficient independent replication.

\subsection{Project Content and Benchmark Status}
EIG measures parameter information through the grid sPCE formulation.
MSC and MOCU instead use the separately defined terminal-control problem,
with numerical safety assessment, posterior coverage and a model-dependent
oracle. These objectives must retain separate result tables and units.
The signed endpoint-RoCoF conference campaign is one explicit protocol;
maximum-absolute-RoCoF and sampled-frequency variants remain distinct.

SIR-ODE is a separate sequential-design benchmark with its own likelihood,
design variables and finite-particle information estimator. Its information
scores must not be relabeled as the power-grid sPCE bound. The optional SIR
MoE uses four experts with top-two sparse execution, whereas the current
direct grid MoE evaluates all four experts and blends latent actions with
soft weights. Architecture, training seeds and estimator settings must
therefore be reported per benchmark rather than inferred across systems.

IEEE9 supports the project's non-reset grid formulation. Reduced IEEE14 is
validated for the declared EIG protocol; control-oriented IEEE14 and the
referenced full dynamic IEEE30 extension require separate validation.
The current conference emphasizes IEEE9/IEEE14 EIG and does not exhaust
the contents or intended extensions of this project.

\subsection{Scope of the Benchmarks}
SIR remains a separate EIG implementation benchmark. Its former results do
not validate the changed power-grid observation or control protocol.
The reduced five-machine IEEE14 EIG backend passed online state-carry and
CPU/GPU validation; its conference campaign uses the signed endpoint sensor.
IEEE14 MSC/MOCU remain outside the validated active scope.
The planned IEEE30 extension retains Demetriou et al.'s dynamic reference
with six machines, including two governed generators and four synchronous
condensers, and six added machine terminal buses~\cite{demetriou2017dynamic}.
Its proposed uncertainty is six inertia constants and two generator governor
parameters, rather than assigning $M/K$ to every electrical bus. The full
machine/exciter/governor implementation and initialization remain pending.
Only completed and verified IEEE14 EIG results may enter the comparison;
no IEEE30 results are claimed.

\subsection{Modeling Assumptions and Interpretation Limits}
The probe sequence now changes both what is known about the grid and the
state from which terminal control acts. This is why the physical oracle is
history dependent and why a state-independent control bank is insufficient.
MSC directly values the selected control under a posterior safety constraint.
MOCU evaluates excess loss under its declared surrogate; EIG values parameter
information. These are different objectives, even where MSC and MOCU select
the same controller for a given history.

Engineering interpretation still requires validation of the reduced dynamics,
parameter priors, contingency, sensing assumptions, control authority and
physical limits. Neither coverage nor numerical agreement with a reduced
simulator certifies field safety. Zero decision latency, a single prescribed
contingency and a scalar step actuator limit the present application.
The MOCU shortfall loss is not a validated outage-damage model. Method claims
must also distinguish original algorithms from the documented objective and
numerical-optimization adaptations. Adaptivity and non-myopic gains remain
empirical questions rather than required outcomes imposed on the experiment.

\section{Results and Discussion}
\label{sec:results}
Historical reset results and the earlier repeated-duration EIG pilot are
separate protocols. They must not populate the new comparison tables. A
successful smoke test establishes workability only, and no method superiority
is concluded from it. Superseded 16-update Step-DAD results are excluded;
the conference uses four-update Step-DAD. MoE-sBOED remains optional and
is not a seventh required row in the primary comparison.
% Populate this section only from completed, validated runs. Do not insert
% smoke scores, incomplete checkpoints, inferred rankings or invented values.
\subsection{Design Performance Across Horizons}
% EIG: methods as rows; T=3,4,5 as columns; report sPCE and contrast count.
% MSC/MOCU: separate metric tables only when matching experiments are complete.
% State the actual training/evaluation seed coverage; a single seed pair has
% no across-seed standard deviation. Never pool incompatible protocols.
\subsection{Optimization Convergence and Numerical Sensitivity}
% Present initialization and validation histories, selected checkpoints,
% late-training changes, and completed particle/contrast/gradient checks.
\subsection{Computational Cost}
% Separate offline training, online decisions/refinement and evaluation-only
% computations. Include hardware and attribute shared Fixed/DAD training costs.
\subsection{Interpretation of the Comparisons}
% Discuss observed adaptivity and lookahead effects only after results exist.
% Include controller safety/oracle diagnostics only for evaluated control tasks.

\section{Conclusion}
\label{sec:conclusion}
The formulation connects an ordered sequence of noisy, non-reset probe
experiments to a continuous terminal power-injection decision. It separates
EIG, minimum safe control and finite-loss MOCU, and specifies the corresponding
physical safety evaluation. \todo{Add conclusions supported by completed
non-reset experiments after the required validation.}

\bibliographystyle{IEEEtran}
\bibliography{objective_driven_boed_refs}
\end{document}
